%──────────────────────────────────────────────────────────────────────────────
%  Rodrigo França Chaves — Curriculum Vitae
%  Compiler: pdflatex
%──────────────────────────────────────────────────────────────────────────────
\documentclass[10pt, letterpaper]{article}

\usepackage[T1]{fontenc}
\usepackage[utf8]{inputenc}

\usepackage{charter}

\usepackage[
  top=0.5in, bottom=0.5in,
  left=0.65in,  right=0.65in
]{geometry}

\usepackage[dvipsnames]{xcolor}
\definecolor{linkblue}{HTML}{1155CC}

\usepackage[
  colorlinks=true,
  urlcolor=linkblue,
  linkcolor=linkblue
]{hyperref}

\usepackage{enumitem}
\setlist[itemize]{
  leftmargin=1.4em,
  itemsep=0pt,
  parsep=0pt,
  topsep=1pt,
  label={\small$\bullet$}
}

\usepackage{tabularx}
\usepackage{titlesec}
\usepackage{array}

% Improves justification so fewer words break across lines
\usepackage{microtype}
\hyphenation{peer-reviewed pro-duced iden-ti-fiers}

\pagestyle{empty}
\setlength{\parindent}{0pt}
\setlength{\parskip}{0pt}

\titleformat{\section}
  {\normalsize\bfseries\uppercase}
  {}{0pt}{}
  [\vspace{0.5pt}\titlerule\vspace{1.5pt}]

\titlespacing*{\section}{0pt}{5pt}{2.5pt}

\newcommand{\cventry}[4]{%
  \noindent
  \begin{tabularx}{\linewidth}{@{}X r@{}}
    \textbf{#1} & \textit{\small #4} \\[0pt]
    \textit{\small #2, #3} & \\
  \end{tabularx}%
  \vspace{0pt}%
}

\newcommand{\cvedu}[4]{%
  \noindent
  \begin{tabularx}{\linewidth}{@{}X r@{}}
    \textbf{#1} & \textit{\small #4} \\[0pt]
    \textit{\small #2, #3} & \\
  \end{tabularx}%
  \vspace{0pt}%
}

\newcommand{\cvpub}[3]{%
  \noindent\textbf{#1}\\[0.5pt]
  {\small\textit{#2}}\\[0.5pt]
  {\small #3}%
  \par\vspace{7pt}%
}

\newcommand{\cvproject}[4]{%
  \noindent\textbf{#1} \enspace{\small\texttt{#2}}\\[0.5pt]
  {\small #4 \enspace #3}%
  \par\vspace{4pt}%
}

\newcommand{\skillrow}[2]{%
  \noindent\textbf{#1:} \enspace #2\par\vspace{1.5pt}%
}



\begin{document}

%── HEADER ───────────────────────────────────────────────────────────────────
\begin{center}
  {\LARGE\bfseries Data Scientist \& Economic Analyst}\\[3pt]
  {\large Rodrigo França Chaves}\\[3pt]
  {\small
    Chicago, IL, USA \enspace$\cdot$\enspace
    \href{mailto:rodrigofrancachaves99@gmail.com}{rodrigofrancachaves99@gmail.com} \enspace$\cdot$\enspace
    \href{mailto:rchaves@uchicago.edu}{rchaves@uchicago.edu} \enspace$\cdot$\enspace
    +1 312 217 5526
  }\\[1.5pt]
  {\small
    \href{https://rodrigofch7.github.io}{Website} \enspace$\cdot$\enspace
    \href{https://scholar.google.com.br/citations?user=6D65dqUAAAAJ}{Google Scholar} \enspace$\cdot$\enspace
    \href{https://www.linkedin.com/in/rodrigofrancac}{LinkedIn} \enspace$\cdot$\enspace
    \href{https://github.com/Rodrigofch7}{GitHub} \enspace$\cdot$\enspace
    Brazilian citizen, F1 visa (STEM OPT eligible)
  }
\end{center}

\vspace{1pt}
{\rule{\linewidth}{0.8pt}}
\vspace{3pt}


{\small
Applied economist and data scientist working on cities: property taxation, housing markets, transit, utility regulation, and crime. First author of a peer-reviewed article in \textit{Utilities Policy} on the health effects of water privatization across two decades of Brazilian municipal data. Builds evidence from sources that resist analysis — assessment rolls, transaction records, disease registries, century-old archives, satellite imagery — pairing quasi-experimental identification with machine learning at scale, and explaining the result to people who did not build the model.
\par}


%── EDUCATION ─────────────────────────────────────────────────────────────────
\section{Education}

\cvedu{M.S. in Computational Analysis and Public Policy}{The University of Chicago}{Chicago, IL}{Sep 2025 – Jun 2027}
{\small Specializations: Markets \& Regulation; Health Policy \\
Coursework: Machine Learning for Public Policy, Applied Econometrics, Data Engineering, Program Evaluation, Public Finance, Statistical Computing}

\vspace{3.5pt}

\cvedu{Postgraduate Degree in Data Science (Pós-Graduação)}{Insper – Institute of Research and Education}{São Paulo, Brazil}{Sep 2022 – Apr 2024}
{\small Coursework: Statistical Inference, Econometrics, Machine Learning, Cloud Computing, Predictive Modeling}

\vspace{3.5pt}

\cvedu{B.A. in International Relations}{Ibmec – Brazilian Institute of Capital Markets}{Belo Horizonte, Brazil}{Feb 2018 – Jun 2022}
{\small Coursework: Political Economy, International Trade, Macroeconomics, Quantitative Methods, Public Policy, Political Science}


%── PUBLICATIONS & WORKING PAPERS ────────────────────────────────────────────
\section{Publications \& Working Papers}

\cvpub{\href{https://www.sciencedirect.com/science/article/abs/pii/S0957178725000189}{Private Ownership of Water and Wastewater Systems: Assessing Health Impacts}}
      {Utilities Policy, peer-reviewed journal article \textnormal{(first author)}}
      {In 2020 Brazil made it easier for private companies to run water and sewage systems. This paper asks whether that was a good idea by looking at the municipalities that had already privatised. Comparing closely matched municipalities before and after they switched, over 1998 to 2021, we find fewer hospital admissions for diarrhoeal disease among children under five (ATT $-11.6$ for infants, $-15.8$ at ages 1 to 5) and no change in diseases unrelated to sanitation. Results vary across municipalities, and the clearest evidence comes from Tocantins, the one state that moved to full private ownership. I wrote the paper, built the dataset, and wrote the code; my co-author supervised.}

\cvpub{\href{https://papers.ssrn.com/sol3/papers.cfm?abstract_id=5441274}{Resistance to Doing Right: Rising Cost of Misconduct from Body-Worn Camera Adoption}}
      {Working paper, in progress \textnormal{(second author)}}
      {Between 2021 and 2022 São Paulo issued 10,125 body cameras across 64 police battalions. We use confidential Military Police records from 2020 to 2023 to ask what changed. Battalion boundaries are not public, so I rebuilt them by clustering the locations of recorded incidents into Voronoi polygons, which is what made it possible to say which areas were treated. Confrontations between police and civilians fell (ATT $-0.50$, significant at 1\%), with the drop concentrated in low-income areas (5\%, and no significant change in high-income ones), while property crime rose, most clearly vehicle theft. That pattern looks less like better policing and more like less policing.}

%── WORK EXPERIENCE ──────────────────────────────────────────────────────────
\section{Work Experience}

\cventry{Data Science and Research Intern, DECDI (formerly DIME)}{World Bank}{Washington, DC}{Jun – Aug 2026}
\begin{itemize}
  \item Measured how construction in Dakar responded to a new bus rapid transit corridor, comparing 100\,m grid cells near the route against cells further out using Google Open Buildings and GHSL satellite data. Because the line was announced in 2018 and only opened in 2024, the design captures how the city reacted in anticipation and during construction, not the effect of the line running.
  \item Found that the project's building-detection threshold had been set at nearly twice the value its source paper recommends for Africa. The old setting discarded real buildings hardest in dense and informal areas, a bias that varies by neighbourhood and so would not have cancelled out of the comparison. Shipped the corrected pipeline with 200 tests and an interactive explorer for the team.
  \item Built a single timeline of forty years of Mumbai property tax and land-use policy, then tested millions of assessment records for bunching at policy thresholds, the pattern that appears when declared values are being manipulated.
  \item Measured which instruments developers actually use to buy extra floor area in Mumbai (TDR, Premium FSI, Fungible FSI) across two decades of transactions, and what market, regulatory, and political conditions drive that choice. \textit{(Findings confidential to the World Bank.)}
\end{itemize}

\vspace{3pt}

\cventry{Graduate Research Assistant}{Mansueto Institute for Urban Innovation, UChicago}{Chicago, IL}{Mar – Jun 2026 · Resuming Sep 2026}
\begin{itemize}
  \item Turned century-old fire insurance registers, cloth-bound books that no computer could read, into usable data. I built the whole pipeline: cutting the pages, running OCR, and cleaning the output into structured records.
  \item Pulled thousands of company addresses out of hundreds of register pages, matched them to three historical censuses by street name, and used NLP on the company names to sort establishments into hundreds of industries. The result is a picture of where industry sat in 20th-century Chicago that nobody had before.
  \item Returning in September 2026 to write the method up as a paper, so other researchers working with archival records can reuse it.
\end{itemize}

\vspace{3pt}

\cventry{Research Assistant}{Insper Cities Lab (Arq.Futuro)}{São Paulo, Brazil}{Aug 2023 – Jul 2025}
\begin{itemize}
  \item Two years evaluating public policy in Brazil across water utilities, crime, property taxation, and real estate markets, using difference-in-differences, regression discontinuity, and spatial methods.
  \item Built the underlying data: home prices, tax assessments, and disease rates, matched across municipalities and census tracts, and grouped into new geographic units through spatial clustering where no official boundary suited the question.
  \item Reconciled administrative records from several government sources, resolving mismatched identifiers and gaps in coverage so the panels lined up across two decades. Without that step the privatisation analysis would not have been identifiable at all.
  \item First author on the resulting \textit{Utilities Policy} article, second author on an in-progress working paper about police body cameras.
  \item Mentored undergraduate researchers on research design, econometrics, and academic writing, and coordinated teams working to publication deadlines.
\end{itemize}

\vspace{3pt}

\cventry{Teaching Assistant: GIS, Econometrics and R}{Insper}{São Paulo, Brazil}{Aug 2023 – Jul 2025}
\begin{itemize}
  \item Taught GIS and R to more than 50 students, using São Paulo as the case study: its housing market, its transit network, and the way both density and prices fall as you move away from the centre.
  \item Advised undergraduate theses through Brazil's required capstone. One placed second in its cohort.
  \item Ran an intensive field course that required working directly with city officials.
\end{itemize}


%── PROJECTS ─────────────────────────────────────────────────────────────────
\section{Projects}

\cvproject{NYC Property Tax: A Horizontal-Equity Audit}
          {Python · LightGBM · scikit-learn}
          {\href{https://github.com/Rodrigofch7/nyc_property_taxes_local}{[GitHub]} \enspace \href{https://nycpropertytaxes.streamlit.app/}{[Live Dashboard]}}
          {Two houses that are alike should pay similar property tax. This project checks whether that holds in New York. Comparing every property against a citywide average is meaningless when a Manhattan co-op and a Queens two-family land in the same column, so each property is instead compared only against properties genuinely like it: same borough, building class, tax class, decade built, size, and price range, with a five-level fallback so rare building types still get a comparison group. On that basis roughly 43\% sit more than 15\% away from what their closest comparables pay. How much you find depends on how a comparable is defined, which was our choice, so the figure measures how spread out assessments are within similar groups rather than proving anyone was assessed wrongly. Narrowing the price bands from quintiles to ventiles cut the price spread inside a group from \$475k to \$48k, which cost a little accuracy and bought labels that actually mean something. Built on 138 leakage-free features and LightGBM, delivered as a dashboard for audit triage. I designed the comparison method, built the data pipeline, and trained the models; two collaborators contributed.}

\cvproject{Data Centers Next Door: Cloud Infrastructure and Housing Costs}
          {R · Shiny · Spatial Econometrics}
          {\href{https://github.com/Rodrigofch7/data-centers-urban-effects}{[GitHub]} \enspace \href{https://rodrigofrancac.shinyapps.io/project-datacenter-urban-effects/}{[Live Dashboard]}}
          {Data centres are going up fast around Chicago, and people living nearby want to know what that does to their bills and their home values. This project puts the question on a map. I scraped 124 listed facilities, confirmed 45 of them, geocoded them, and worked out roughly when each opened from air permit filings between 2000 and 2024. Those dates are lined up against ZIP-level Zillow home values and Census household and utility costs across 660 metro ZIP codes, so a user can see what moved before and after each opening. It shows patterns rather than causes: a facility opening near rising prices does not mean it caused them. Course project with four people; I did the scraping, the geocoding, and the Shiny dashboard.}


%── TECHNICAL SKILLS ─────────────────────────────────────────────────────────
\section{Technical Skills}

\skillrow{Methods}{Causal Inference (DiD, Staggered DiD, RDD, IV), Panel \& Spatial Econometrics, Remote Sensing}
\skillrow{Programming \& Data}{Python, R, SQL, PostgreSQL, Stata; data pipelines, OCR, geospatial workflows}
\skillrow{Tools}{QGIS, ArcGIS, Git, LaTeX, Excel, Streamlit, Shiny}
\skillrow{Domain}{Property tax \& valuation, zoning \& land use, housing \& infrastructure policy, utility regulation}
\skillrow{Languages}{Portuguese (Native), English (Fluent, TOEFL 112/120), Spanish (Conversational), German (Basic)}


\end{document}
